% !TEX root = ../memorianueva.tex
% Chapter 2

\chapter{Introducción específica} % Main chapter title

\label{Chapter2} % Para referenciar este capítulo: \ref{Chapter2}

En el presente capítulo se describen las tecnologías, plataformas, servicios y componentes de terceros utilizados como base para el desarrollo del sistema. Se presentan los dispositivos de borde, los protocolos de comunicación, las plataformas en la nube, los servicios de inteligencia artificial, las herramientas de almacenamiento y los recursos de visualización empleados en la solución.

%----------------------------------------------------------------------------------------
%	SECTION 1
%----------------------------------------------------------------------------------------

\section{Dispositivos de borde}
\label{sec:dispositivos_borde}

El nodo de borde del sistema está constituido por un servidor local instalado en el comercio y por una impresora térmica conectada a él. A diferencia de un sistema de Internet de las Cosas tradicional, no se emplean microcontroladores ni sensores. El dispositivo de borde es un equipo de cómputo de propósito general que ejecuta los servicios locales y comanda la impresión de los pedidos.

\subsection{Servidor local}
\label{subsec:servidor_local}

Como servidor local se utiliza una computadora con sistema operativo Windows Server 2019 \cite{windowsserver}, propiedad del cliente y preexistente en el local. Sobre este equipo se ejecutan el servidor de impresión y el cliente del túnel seguro hacia la nube.
% [PENDIENTE: si se conocen, agregar las características básicas del equipo
% (procesador, memoria).]

Se evaluaron dos alternativas para el nodo de borde. La primera consistió en incorporar una computadora de placa única (SBC, del inglés \textit{Single Board Computer}) dedicada, como una Raspberry Pi \cite{raspberrypi}. La segunda consistió en reutilizar el equipo ya disponible en el comercio. Se optó por la segunda opción por tres motivos. El primero fue el costo, dado que reutilizar el hardware disponible elimina la inversión marginal en equipamiento. El segundo fue la operación, ya que el equipo permanece encendido durante el horario comercial y resulta conocido por el personal, lo que simplifica la instalación y el mantenimiento. El tercero fue la compatibilidad, porque la impresora provee controladores para Windows y su integración no requirió desarrollos de bajo nivel.

Como contrapartida, un equipo de propósito general consume más energía que una SBC dedicada y puede verse condicionado por otros usos dentro del comercio. Esta limitación se consideró aceptable frente a la simplicidad operativa obtenida.

También se descartó el uso de un microcontrolador, ya que el conjunto de responsabilidades del nodo de borde (un servidor web local, el cliente del túnel seguro y la administración de la cola de impresión) se resuelve de manera más conveniente sobre un sistema operativo completo.

\subsection{Impresora térmica}
\label{subsec:impresora}

Para la impresión de los tickets se utiliza una impresora térmica Hasar P-HAS-181 \cite{hasar181}, un modelo de tipo comandera del fabricante argentino Grupo Hasar, orientado al punto de venta. La impresora imprime por cabezal térmico directo sobre rollos de papel de 80 mm de ancho y dispone de tres interfaces de conexión (puerto serie, USB y Ethernet) y de controladores para Windows. Sus comandos de impresión resultan compatibles con ESC/POS \cite{escpos}, el conjunto de instrucciones definido por Epson que constituye el estándar de facto del sector.

En la instalación actual, la impresora se conecta al servidor local por USB. La disponibilidad de una interfaz Ethernet fue, sin embargo, un criterio de selección. El equipo puede requerir un cambio de ubicación dentro del comercio, por ejemplo hacia el sector de armado de pedidos. La conexión por red permitiría ese traslado sin dependencia de la proximidad física con el servidor.

La elección de este modelo se fundamentó en cuatro factores. La impresión térmica es el estándar del comercio minorista por su velocidad y su bajo costo operativo, al no utilizar tinta ni cartuchos. La compatibilidad con ESC/POS permite comandar el equipo desde bibliotecas ampliamente disponibles, sin depender de un protocolo propietario. Los controladores para Windows garantizan la integración directa con el servidor local. Por último, se trata de un fabricante nacional, lo que asegura la provisión local de insumos, repuestos y servicio técnico, un criterio relevante para un comercio sin personal técnico propio.

%----------------------------------------------------------------------------------------
%	SECTION 2
%----------------------------------------------------------------------------------------

\section{Protocolos de comunicación}
\label{sec:protocolos}

El intercambio de mensajes entre los componentes se realiza mediante HTTP \cite{rfc9110}, con interfaces que siguen el estilo arquitectónico REST \cite{fielding2000rest}. La elección responde a que todos los servicios involucrados exponen sus funcionalidades a través de interfaces de programación de aplicaciones (API, del inglés \textit{Application Programming Interface}) basadas en HTTP. Para los eventos asincrónicos, como la llegada de un mensaje, se emplean notificaciones automáticas (\textit{webhooks}), que evitan que el sistema receptor deba consultar de forma periódica por novedades.

\subsection{Canal de mensajería}
\label{subsec:whatsapp}

La vinculación con WhatsApp se realiza mediante Evolution API \cite{evolutionapi}, una pasarela de mensajería de código abierto que expone las funciones de la plataforma a través de una interfaz REST y de notificaciones automáticas. Evolution API admite dos modos de conexión. El primero se basa en el protocolo de WhatsApp Web, mediante la biblioteca Baileys \cite{baileys}, y no requiere verificación comercial ante Meta. El segundo se basa en la API oficial WhatsApp Business Cloud \cite{whatsappbusiness}. En este trabajo se utiliza el primer modo. Ante cada mensaje entrante, de texto, voz o imagen, la pasarela emite una notificación con el contenido o con las referencias para descargarlo.

La elección de Evolution API para la etapa inicial del trabajo se fundamentó en su costo nulo, en su rapidez de puesta en marcha (no requiere verificación comercial ni aprobación de plantillas) y en su despliegue autoalojado, coherente con la infraestructura del sistema. Su principal limitación es que el modo basado en WhatsApp Web constituye un mecanismo no oficial, con posibles restricciones de estabilidad y cumplimiento. Por ese motivo, el sistema se diseñó compatible con la API oficial de WhatsApp Business, y la migración prevista no implica modificaciones estructurales en los flujos de procesamiento.

\subsection{Comunicación segura con la nube}
\label{subsec:cloudflare}

Para vincular los servicios en la nube con el servidor local sin apertura de puertos en la red del comercio, se emplea un túnel seguro provisto por Cloudflare \cite{cloudflaretunnel}. Un cliente ligero, ejecutado en el servidor local, establece una conexión cifrada saliente hacia la infraestructura del proveedor. A través de esa conexión, los servicios remotos alcanzan al nodo de borde.

Se consideraron dos alternativas. La apertura de puertos en el enrutador del local exige configurar el equipamiento de red y deja un servicio a la escucha de conexiones entrantes desde Internet, lo que amplía la superficie de ataque. Una red privada virtual (VPN, del inglés \textit{Virtual Private Network}) resuelve la seguridad, pero agrega administración de claves y configuración en ambos extremos. El túnel de Cloudflare, en cambio, no requiere configuración del enrutador, ya que la conexión es saliente, cifra el tráfico de extremo a extremo y mantiene la red interna sin exposición entrante, con una operación mínima acorde a un local sin personal técnico.

%----------------------------------------------------------------------------------------
%	SECTION 3
%----------------------------------------------------------------------------------------

\section{Plataformas cloud y servicios}
\label{sec:plataformas_cloud}

En esta sección se describen los servicios en la nube sobre los que se apoya el sistema, que comprenden la plataforma de orquestación de los flujos de procesamiento, la infraestructura donde se despliegan los componentes y los servicios de inteligencia artificial que resuelven la interpretación de los pedidos.

\subsection{Orquestación de flujos de trabajo}
\label{subsec:n8n}

La orquestación de las tareas de recepción, interpretación y estructuración de los pedidos se realiza con n8n \cite{n8n}, una plataforma de automatización de flujos de trabajo de código abierto. En n8n, la lógica se construye mediante la conexión de nodos configurables (disparadores, transformaciones e integraciones con servicios externos) en flujos visuales. Este enfoque permite integrar servicios heterogéneos sin desarrollar desde cero la lógica de integración. La plataforma incorpora además agentes de inteligencia artificial (\textit{AI Agents}), que encadenan modelos de lenguaje con herramientas y fuentes de datos. Sobre esa capacidad se apoya la interpretación de los pedidos.

Frente a alternativas comerciales de automatización en la nube, como Zapier \cite{zapier}, n8n presentó dos ventajas decisivas para este trabajo. La primera es la posibilidad de autoalojamiento en la propia infraestructura, sin costo por tarea ejecutada, un factor crítico dado el volumen diario de mensajes. La segunda es que los datos de los clientes permanecen dentro del entorno controlado por el sistema, sin atravesar una plataforma de terceros adicional.

\subsection{Infraestructura de despliegue}
\label{subsec:vps_docker}

Los servicios en la nube se alojan en un servidor privado virtual (VPS, del inglés \textit{Virtual Private Server}) del proveedor Hostinger \cite{hostinger}, en su plan KVM 4. Este plan provee 4 núcleos virtuales, 16 GB de memoria RAM, 200 GB de almacenamiento NVMe y 16 TB de transferencia mensual. El dimensionamiento respondió a la carga concurrente del sistema, ya que sobre el mismo servidor conviven la pasarela de mensajería, el motor de flujos, las bases de datos y el panel web.

Todos los servicios se despliegan de forma contenerizada mediante Docker \cite{docker}. La contenerización empaqueta cada servicio junto con sus dependencias en unidades aisladas y reproducibles, lo que permite describir el despliegue en archivos versionables, reproducir el entorno sin diferencias de configuración y actualizar cada componente de manera independiente.

\subsection{Servicios de inteligencia artificial}
\label{subsec:servicios_ia}

La interpretación de los pedidos requiere resolver tres tareas diferenciadas: convertir los mensajes de voz en texto, leer las listas manuscritas enviadas como fotografía y extraer de un texto informal los productos, las cantidades y las unidades solicitadas. Cada tarea se resuelve con un modelo de inteligencia artificial consumido como servicio en la nube, provisto por OpenAI \cite{openai} o por Anthropic \cite{claude}.

La transcripción de audio se realiza con el servicio de transcripción de OpenAI, basado en el modelo Whisper \cite{whisper}. Su característica más relevante para este trabajo es la robustez frente a condiciones adversas de grabación, como el ruido de fondo, la variabilidad de acentos y la calidad desigual del micrófono. Los mensajes de voz de los clientes se graban dentro del comercio o en la vía pública, sin control sobre el entorno sonoro, por lo que esa robustez resultó determinante.

La lectura de las imágenes se resuelve con modelos multimodales de visión. Se evaluó como alternativa un motor de reconocimiento óptico de caracteres (OCR, del inglés \textit{Optical Character Recognition}) clásico, pero las imágenes recibidas (listas manuscritas con caligrafías dispares e iluminación irregular, y capturas de planillas) le hacen perder precisión, porque opera sobre cada carácter de manera aislada y sin considerar la disposición del contenido. Los modelos multimodales, en cambio, interpretan la imagen en su contexto, reconocen la estructura tabular y aprovechan el conocimiento del dominio para resolver ambigüedades.

Sobre cada imagen se ejecutan dos modelos de proveedores distintos, GPT-4o de OpenAI y Claude Sonnet de Anthropic, porque ante una misma imagen producen en ocasiones lecturas discrepantes, sobre todo con caligrafías difíciles o planillas complejas. Las dos lecturas se reconcilian en un tercer paso, donde el modelo GPT-5.5 de OpenAI las compara, corrige las diferencias y devuelve una transcripción única. Esta redundancia agrega costo y latencia, pero reduce el riesgo de un error silencioso en la etapa más frágil de la interpretación.

La estructuración final se resuelve con un modelo extenso de lenguaje (LLM, del inglés \textit{Large Language Model}), GPT-5.5 de OpenAI, integrado como agente dentro del flujo. Recibe el texto del pedido junto con el catálogo del comercio y produce una representación estructurada con los productos, las cantidades y las unidades. Su capacidad de interpretar lenguaje informal, con abreviaturas, errores de tipeo y denominaciones no normalizadas, constituye el aporte central de estos modelos a la solución.

La decisión de consumir estos modelos como servicio, en lugar de ejecutarlos sobre la propia infraestructura, define la distribución del cómputo entre el borde y la nube. Los tres modelos exigen recursos de procesamiento muy superiores a los del nodo de borde, por lo que su ejecución local no resultaba viable con el equipamiento del comercio. Su consumo como servicio mantiene la inversión inicial y la complejidad operativa en el mínimo. El costo es doble: un componente variable por uso y la dependencia de la conectividad para la etapa de interpretación. La actuación crítica del sistema permanece en el borde, de manera que esa dependencia no compromete la operación local de impresión.

%----------------------------------------------------------------------------------------
%	SECTION 4
%----------------------------------------------------------------------------------------

\section{Herramientas de almacenamiento y visualización}
\label{sec:almacenamiento_visualizacion}

En esta sección se presentan las herramientas de persistencia de datos y los componentes destinados a la visualización de la operación. Las bases de datos que sostienen la información del sistema y el servidor local junto con el panel web de control.

\subsection{Bases de datos}
\label{subsec:bases_datos}

La información de pedidos, clientes y productos se almacena en PostgreSQL \cite{postgresql}, una base de datos relacional de código abierto. Su modelo relacional con integridad referencial resulta adecuado para el dominio del problema, donde los pedidos se vinculan con clientes y con artículos del catálogo. La trazabilidad exigida al sistema requiere garantías de consistencia transaccional, y las consultas sobre el historial alimentan además las métricas del panel de control.

De manera complementaria se emplea Redis \cite{redis}, una base de datos en memoria de tipo clave-valor, para la gestión de las sesiones y del estado de los clientes activos durante la recepción de un pedido. Redis ofrece tiempos de acceso muy bajos y expiración automática de claves, características apropiadas para datos de carácter transitorio con alta frecuencia de lectura y escritura. La combinación responde a un criterio de diseño explícito: los datos persistentes y auditables residen solo en la base relacional, mientras que Redis se reserva para información temporal, de modo que una eventual falla de la capa en memoria no compromete los registros del negocio.

\subsection{Servidor local y panel de visualización}
\label{subsec:flask_panel}

El servidor de impresión que se ejecuta en el nodo de borde se implementó sobre Flask \cite{flask}, un entorno de desarrollo web minimalista escrito en el lenguaje Python. Flask permite exponer servicios HTTP con una huella de recursos reducida y sin la complejidad de un servidor de aplicaciones completo. Resulta adecuado para un servicio local de propósito específico, que recibe el pedido estructurado y comanda la impresora sobre el equipo Windows del comercio.

La visualización de los pedidos, del historial y de las métricas de operación se ofrece a través de un panel web alojado en el VPS, accesible desde cualquier navegador moderno. Este enfoque evita la instalación de aplicaciones adicionales en los dispositivos del comercio. El panel se desarrolló con Angular \cite{angular}, un entorno de desarrollo de aplicaciones web mantenido por Google, junto con la biblioteca de componentes Angular Material. La representación de las métricas se resuelve con Chart.js \cite{chartjs}. La elección de un entorno estructurado, con tipado estático y organización por componentes, respondió a la necesidad de sostener el crecimiento del panel a medida que se incorporan nuevas vistas y funciones de administración.